\icode{N24G-MCH-SI-L3-261026-002}
\nceeID[6]{2024江西}{8}
一条河流某处存在高度差，小鱼从低处向上跃出水面，冲到高处。如图所示，以小鱼跃出水面处为坐标原点，\(x\) 轴沿水平方向，建立坐标系，小鱼的初速度为 \(v_0\)，末速度 \(v\) 沿 \(x\) 轴正方向。在此过程中，小鱼可视为质点且只受重力作用。关于小鱼的水平位置 \(x\)、竖直位置 \(y\)、水平方向分速度 \(v_x\) 和竖直方向分速度 \(v_y\) 与时间 \(t\) 的关系，下列图像可能正确的是\paren[]

\begin{tikzpicture}[>=stealth,font=\scriptsize]
	\draw[->,blue] (-.3,0)--(4,0)node[below]{ \(x\)};
	\draw[->,blue] (0,0)node[below]{ \(O\)}--(0,3)node[left]{ \(y\)};
	\draw[dashed,cyan] (2.8,2.2) parabola bend(2.8,2.2)(0,0);
	\draw (2.1,0)--(2.7,2)--(3.5,2);
	\draw[fill=white] (2.8,2.2)..controls(2.95,2.1)..(3.1,2.2)..controls(2.95,2.3)..(2.8,2.2)--++(150:.1)--++(0,{-.2*sin(30)})--(2.8,2.2);
	\fill (3,2.2) circle(.03);
	\draw[->,red] (0,0)--++(60:1)node[above left]{ \(v_0\)};
\end{tikzpicture}

\begin{tikzpicture}[>=stealth,font=\scriptsize]
	\begin{scope}
		\draw[->,blue] (0,0)node[below left]{ \(O\)}--(2.6,0)node[below]{ \(t\)};
		\draw[->,blue] (0,0)--(0,2.6)node[left]{ \(x\)};
		\draw[magenta] (0,0)--(2,1.8);
		\node at(1.3,-.4){\normalsize A};
	\end{scope}
	\begin{scope}[xshift=3.5cm]
		\draw[->,blue] (0,0)node[below left]{ \(O\)}--(2.6,0)node[below]{ \(t\)};
		\draw[->,blue] (0,0)--(0,2.6)node[left]{ \(y\)};
		\draw[magenta] (0,0) parabola bend(0,0) (2,2.2);
		\node at(1.3,-.4){\normalsize B};
	\end{scope}
	\begin{scope}[xshift=7cm]
		\draw[->,blue] (0,0)node[below left]{ \(O\)}--(2.6,0)node[below]{ \(t\)};
		\draw[->,blue] (0,0)--(0,2.6)node[left]{ \(v_x\)};
		\draw[magenta] (0,.4)--(2,2.2);
		\node at(1.3,-.4){\normalsize C};
	\end{scope}
	\begin{scope}[xshift=10.5cm]
		\draw[->,blue] (0,0)node[below left]{ \(O\)}--(2.6,0)node[below]{ \(t\)};
		\draw[->,blue] (0,0)--(0,2.6)node[left]{ \(v_y\)};
		\draw[magenta] (0,2.2)--(2,0);
		\node at(1.3,-.4){\normalsize D};
	\end{scope}
\end{tikzpicture}
\begin{answer}
	A\quad 提示：\(x\)方向匀速运动，\(y\)方向匀减速运动
\end{answer}